\chapter{Documentation and the Demo That Gets You Called Back}
\label{ch:projects-c}

\section{Documentation That Sells Your Skills}

\begin{keytakeaway}
Your README is your elevator pitch to recruiters. In two minutes of reading it should answer: What problem does
this solve? Why should I care? What makes this impressive? Can you actually build production systems?
\end{keytakeaway}

\subsection{README Structure}

\textbf{Section 1 --- Project overview (30 seconds to hook).} One sentence on what it is; the business problem and
impact; headline result numbers.

\begin{example}
``Restaurant demand forecasting system that reduces food waste by 25\% through ML-powered daily predictions.
Predicts demand for 50+ menu items with 15\% MAPE, deployed as automated daily reports used by 3 local
restaurants. Handles seasonal patterns, weather impact, and special events using gradient boosting and custom
feature engineering.''

This tells a recruiter immediately: a real problem was solved, there's actual business impact, it's deployed and
used, and the technical sophistication is appropriate.
\end{example}

\textbf{Section 2 --- Problem definition (1 minute to understand context).} Business context (who has the
problem, what it costs); current solution and what's wrong with it; your proposed solution and approach; success
metrics and what improvement actually matters.

\textbf{Section 3 --- Technical approach (2 minutes to show depth), 2--3 paragraphs per sub-area.} Data pipeline
(sources, quality challenges, preprocessing); feature engineering (what and why, engineering decisions,
importance results); model development (baseline, models compared, final choice and why, tuning approach);
deployment (architecture, serving approach, infrastructure); monitoring (metrics tracked, drift approach,
retraining strategy).

\textbf{Section 4 --- Results (1 minute to prove value).} Model performance on the test set vs.\ baseline, with
error-analysis highlights; business impact (quantified value, adoption, cost/benefit); technical learnings
(challenges overcome, surprising findings, what you'd do differently).

\begin{example}
``Final model achieves 15\% MAPE, beating the 30\% MAPE baseline by 50\%. Predictions are most accurate for
high-volume items (10\% MAPE) but struggle with new menu items (25\% MAPE). Weather features provided an 8\%
improvement; day-of-week patterns were critical for accuracy. Deployed to 3 restaurants, averaging 25\% waste
reduction (\$12K/month savings per restaurant). The system runs daily at 6\,a.m., generating predictions by
7\,a.m.\ for 10\,a.m.\ ordering. Key learning: the model benefits more from a region-specific holiday calendar
than from generic weather data.''
\end{example}

\textbf{Section 5 --- How to run (30 seconds for easy testing).} Prerequisites, installation, configuration, a
quick-start example. If a recruiter can't run it in five minutes, they won't try. \textbf{Section 6 --- Project
structure.} A directory tree with explanations, showing organization and modularity. \textbf{Section 7 --- Future
improvements.} What you'd add with more time (better features, alternative models, deployment enhancements,
scaling) --- this demonstrates you know the project isn't perfect, you understand trade-offs, and you can
prioritize.

\subsection{Additional Documentation}

Beyond the README: \textbf{API documentation} (endpoints, schemas, example calls, error codes, auth);
\textbf{deployment guide} (infrastructure requirements, step-by-step deployment, configuration,
troubleshooting); \textbf{model card} (intended use, training data, performance metrics, ethical considerations,
limitations and biases --- the industry-standard format); \textbf{data documentation} (sources, collection
methodology, quality issues and handling, privacy/ethics, update frequency); \textbf{experiment log} (hypotheses
tested, results summary, key learnings, decision rationale).

\subsection{Code Quality Standards}

Your code is evaluated too: logical organization and clear naming; docstrings, comments for complex logic, type
hints, a README per module; unit and integration tests that actually pass; graceful error handling with
informative messages and logging; no hardcoded values (environment variables for secrets, config files for
parameters); minimal, pinned dependencies with a clear requirements file; meaningful, logical commit history with
no sensitive data committed. This signals maintainable code, awareness that others will read it, and solid
software-engineering fundamentals --- critical for technical interviews.

\subsection{Visual Documentation}

Diagrams that help: system architecture (how components connect), data flow (how data moves through the
pipeline), model architecture (if complex), and results visualizations (performance over time, feature
importance, error analysis, business impact). Tools: draw.io for diagrams, Matplotlib/Seaborn for plots,
Streamlit/Gradio for interactive demos. A picture is worth a thousand words in a README, especially for complex
systems.

\section{The Demo That Gets You Called Back}

\begin{keytakeaway}
Recruiters won't clone your repo and run setup scripts. They need a live demo they can test in a browser, two
minutes to understand what you built, evidence of technical depth, and proof it actually works.
\end{keytakeaway}

\subsection{Interactive Demo Options}

\begin{longtable}{p{0.22\textwidth} p{0.72\textwidth}}
\toprule
\textbf{Option} & \textbf{Best for} \\
\midrule
\endhead
Web app (Streamlit) & Best for most projects --- quick to build, handles UI automatically, free hosting, perfect
for ML demos. Include: an input form, prediction output with confidence, feature-importance view, performance
metrics, an about page. \\
API + Swagger docs & Technical audiences --- deploy to Heroku/Railway with auto-generated, interactive API docs.
Shows production-ready deployment and API design. \\
Jupyter Notebook + Voil\`a & Exploration-heavy projects --- converts a notebook to a web app with interactive
widgets, no code visible. Use when visualizations matter more than serving predictions. \\
Video demo (backup) & When a live demo isn't possible --- a 2--3 minute screen walkthrough showing key
functionality and explaining decisions. Not as good as interactive, but better than nothing. \\
\bottomrule
\end{longtable}

\subsection{Demo Design Principles}

Make it obvious what to do: clear instructions (``Enter restaurant capacity and date to get a demand forecast,''
not ``Input data and click submit''); pre-filled example inputs; visual feedback (loading indicators,
success/error states, clear results); input validation that prevents confusion. Show technical depth without
overwhelming: a simple interface for predictions, expandable sections for technical details, optional deep dives.
Highlight key features and confidence scores. Demonstrate robustness by handling edge cases gracefully with
informative messages. Tell the story --- problem, impact, technical approach --- briefly, in the demo itself.

\subsection{What to Show in the Demo}

\textbf{Essential:} prediction capability working reliably; explanation of why (feature contributions,
confidence); performance metrics (test-set results); data insights (feature importance, distributions).
\textbf{Nice-to-have:} comparison to baseline; interactive error-analysis exploration; a simulated monitoring
dashboard; an A/B-testing comparison framework.

\subsection{Deployment Platforms}

\begin{longtable}{p{0.28\textwidth} p{0.64\textwidth}}
\toprule
\textbf{Platform} & \textbf{Good for} \\
\midrule
\endhead
Streamlit Cloud & Free tier, direct GitHub integration, auto-updates on push, custom URLs --- best for interactive
predictions \\
Heroku / Railway & Free/low-cost tiers for APIs, easy deployment, databases available, auto-scaling \\
Hugging Face Spaces & ML demos specifically, free GPU tier, strong for CV/NLP, built-in community visibility \\
Google Colab & Free compute, shareable link, interactive notebooks \\
GitHub Pages & Free static hosting with a custom domain --- good for documentation sites \\
\bottomrule
\end{longtable}

\subsection{Demo Checklist}

\begin{multicols}{2}
\begin{itemize}[nosep]
  \item Works on mobile
  \item Loads in under 5 seconds
  \item Instructions clear on first use
  \item No crashes on weird inputs
  \item Error messages are actually helpful
  \item Professional appearance, no placeholder text
  \item Linked from resume/LinkedIn/GitHub
  \item Uptime checked before interviews
  \item A pre-filled example is provided
  \item An about section gives context
\end{itemize}
\end{multicols}

\begin{keytakeaway}
A working demo is worth roughly ten times more than perfect code nobody can run.
\end{keytakeaway}

\section{Time Investment, Realistically}

\begin{longtable}{p{0.5\textwidth} p{0.42\textwidth}}
\toprule
\textbf{Template} & \textbf{Time investment} \\
\midrule
\endhead
Template 1: End-to-end prediction pipeline (full lifecycle) & 3--4 weeks \\
Template 2: Real-time serving API (low-latency serving) & 2--3 weeks \\
Template 3: Domain-specific application (industry specialization) & 3--4 weeks \\
\bottomrule
\end{longtable}

\noindent (Full detail on these three templates, the 15 fatal mistakes that get projects rejected, and the
interview questions a good project answers for you are consolidated in Appendix~\ref{app:mistakes}--\ref{app:interview} for quick reference.)

\begin{keytakeaway}
The formula: Real Problem + Messy Data + Full Pipeline + Actual Deployment + Clear Documentation = Interview
Calls. Don't try to build three projects in a month --- build one properly in 3--4 weeks (Week 1: problem
definition, data, exploration; Week 2: feature engineering, baseline, experimentation; Week 3: model development,
validation, interpretation; Week 4: deployment, monitoring, documentation). Then build another. By the third
project, you will know more about production ML than most people with the job title.
\end{keytakeaway}

